\section{Geophysical implications}

\subsection{Joint processing of gravity-gradient data}

The common approach to jointly process gravity-gradient data 
\cite[e.g.,][]{boggs-dransfield2004,barnes_lumley2011, martinez_li_2016} or to convert gravity disturbance data into the 
gravity-gradient tensor \cite[e.g.,][]{piauilino_etal2019_seg, oliveirajr_etal2023}, for example, 
usually presume, without any proof, the existence of a single layer of 
monopoles that satisfactorily reproduces the gravitational field and
preserves the physical relationship with the true sources.
This implies that the gravitational potential and its spatial derivatives attributed to the layer are presumed to 
accurately reproduce the corresponding field components produced by the true sources.
We provide the theoretical background for these potential-field transformations.

Let us consider, for example, the common two-step approach to convert gravity disturbance into gravity
gradient data (Table \partoneref{tbl:Table2}). The first step consists in obtaining a particular equivalent layer $p_{\varepsilon}$ 
by considering that the observed potential field $d^{o}$ is given by $\tfrac{\partial}{\partial z} U(\mathbf{r})$ 
(equation \partoneref{eq:potential-U-1st-deriv-space}) and that the kernel function 
$a(\mathbf{r}, \mathbf{r}'')$ is the vertical derivative of $\Psi$ (equation \partoneref{eq:Psi}).
The second step consists in using $p_{\varepsilon}$ to compute the gradient components.
The $uv$-component of the gradient tensor is computed by setting $D_{f} \equiv \tfrac{\partial^{2}}{\partial u \, \partial v}$
(equations \partoneref{eq:potential-U-2nd-deriv-space} and \partoneref{eq:transformed-field-continuous-layer-Neumann}),
where $u, v \in \{ x, y, z \}$.

Note that the first step consists in defining an equivalent layer 
by setting $\alpha_{d} = \alpha_{a} \equiv \alpha(1, \hat{\mathbf{z}})$ in equation \partoneref{eq:p-Fourier}.
According to Definition \partoneref{def:eql-Neumann}, this particular $p_{\varepsilon}$ is a Neumann's equivalent layer,
which naturally preserves the physical relationship with the true sources.
The second step of the common approach corresponds to 
compute transformed fields $f_{\varepsilon}$ (equation \partoneref{eq:transformed-field-continuous-layer-Fourier}) 
defined in the Fourier domain by using $\alpha_{t} = \alpha_{f}$ (equation \partoneref{eq:alpha-t-generic}), with 
$\alpha_{f} \equiv \alpha(1, \hat{\mathbf{u}}) \, \alpha(1, \hat{\mathbf{v}})$ (equation \partoneref{eq:diff-operators-generic}), 
where the unit vectors $\hat{\mathbf{u}}$ and $\hat{\mathbf{v}}$ are given by the canonical base vectors 
(equation \partoneref{eq:canonical-base-vector}).
We can easily get from \partoneref{eq:diff-operators-generic} that this Fourier-domain operator $\alpha_{f}$ is equivalent to the 
space-domain operator $D_{f} \equiv \tfrac{\partial^{2}}{\partial u \, \partial v}$.
Thus, the common two-step approach to convert gravity disturbance into gravity
gradient data via equivalent-layer technique can be theoretically explained as an application of 
a particular Neumann's equivalent layer (Definition \partoneref{def:eql-Neumann}).


\subsection{The existence of a planar dipole layer}


A predominant trend in previous studies utilizing a planar equivalent layer of dipoles for magnetic-field transformations 
is to treat the layer as a purely mathematical abstraction employed for processing total-field anomaly data (Table \partoneref{tbl:Table2}), 
a perspective exemplified by the remarks of \citet{pedersen1989} concerning the research of 
\citet{bhattacharyya_chan1977}.
Specifically, Pedersen pointed out that the first fundamental existence theorem for the Dirichlet's 
problem does not postulate that a general solution can be written as a layer of dipoles magnetized along its 
normal direction. In essence, Pedersen underscored that a dipole layer lacks physical significance.
In response, Chan affirmed that such a layer was employed in his work 
solely to address an integral equation, i.e., as a purely mathematical solution.

Since its inception in the 1960s, the equivalent-layer technique in applied geophysics has garnered 
relatively scant attention in terms of exploring the theoretical attributes of planar dipole layers 
capable of faithfully reproducing total-field anomaly data. 
This inquiry usually aims to discern the physical significance and possible link of such layers 
with boundary problems of potential theory.

Here, we leverage a theoretical approach to infer within the purview of applied geophysics, 
for the first time, the existence of planar dipole layers capable of accurately and simultaneously reproducing 
the magnetic scalar potential and induction field components along any direction, including that representing 
the approximated total-field anomaly (Table \partoneref{tbl:Table2}).

We first consider that (i) the observed potential field $d^{o}$ 
is satisfactorily described by the directional derivative $\hat{\mathbf{u}}^{\top} \mathbf{H}_{U}(\mathbf{r}) \, \hat{\mathbf{v}}$ 
(equation \partoneref{eq:potential-U-oblique-2nd-deriv-space}), where $\hat{\mathbf{u}}$ defines the constant direction of the main geomagnetic
field at the study area and $\hat{\mathbf{v}}$ the constant total-magnetization direction of the sources  (Table \partoneref{tbl:Table2}), and (ii)
the kernel is $a(\mathbf{r}, \mathbf{r}'') = \hat{\mathbf{u}}^{\top} \mathbf{H}_{\Psi}(\mathbf{r}, \mathbf{r}'') \, \hat{\mathbf{h}}$,
where $\mathbf{H}_{\Psi}$ is the Hessian matrix \citep[e.g.,][ p. 158]{kaplan2002} of
function $\Psi$ (equation \partoneref{eq:Psi}) and $\hat{\mathbf{h}}$ is a constant unit vector not necessarily equal to 
$\hat{\mathbf{v}}$.
In this case, we use $\alpha_{d} \equiv \alpha(1, \hat{\mathbf{u}}) \, \alpha(1, \hat{\mathbf{v}})$ and 
$\alpha_{a} \equiv \alpha(1, \hat{\mathbf{u}}) \, \alpha(1, \hat{\mathbf{h}})$ 
(equation \partoneref{eq:diff-operators-generic}) in equation \partoneref{eq:p-Fourier} to obtain the following equivalent layer 
in the Fourier domain:
\begin{equation}
	\mathcal{F} [ p_{\varepsilon}(\mathbf{r}'') ] =
	k \, \tfrac{\left(\mathbf{k}^{\top} \hat{\mathbf{v}} \right)}{\left(\mathbf{k}^{\top} \hat{\mathbf{h}} \right)} \,
	\mathcal{F} \left[ U_{\varepsilon}(\mathbf{r}'') \right] \: .
	\label{eq:p-Fourier-dipoles}
\end{equation}
This particular equivalent layer represents a planar layer of dipoles with constant magnetization direction
$\hat{\mathbf{h}}$ and was deduced by considering
that $\hat{\mathbf{u}}$ is a unit vector having the same constant direction of the main geomagnetic field at the study area.
However, it can be easily verified that the same equivalent layer can also be deduced by considering that $\hat{\mathbf{u}}$ is one of the
canonical base vectors $\hat{\mathbf{x}}$, $\hat{\mathbf{y}}$ or $\hat{\mathbf{z}}$ (equation \partoneref{eq:canonical-base-vector}).
In this case, the observed potential field represents a Cartesian component of the magnetic induction field  (Table \partoneref{tbl:Table2}).
It can also be easily verified that the same $p_{\varepsilon}$ given by equation 
\ref{eq:p-Fourier-dipoles} can be equally deduced if we consider that the 
observed potential field is the magnetic scalar potential  (Table \partoneref{tbl:Table2}). 
In this case, $\alpha_{d}$ and $\alpha_{a}$ are defined
without the term $\alpha(1, \hat{\mathbf{u}})$. 

The main conclusion obtained from the reasoning presented above is that a planar equivalent layer of dipoles defined according 
to equation \ref{eq:p-Fourier-dipoles} simultaneously 
reproduces the magnetic scalar potential and any magnetic induction field component (including the approximated total-field anomaly) 
due to sources with constant total-magnetization direction defined by $\hat{\mathbf{v}}$, even if the layer has a constant magnetization 
direction defined by a unit vector $\hat{\mathbf{h}} \ne \hat{\mathbf{v}}$. 
This theoretical result explains, for example, why we can use a planar equivalent layer of dipoles with arbitrary constant 
magnetization direction $\hat{\mathbf{h}}$ to compute the magnetic induction field components or the amplitude of the magnetic induction field due to the true sources 
without presuming any information about the constant direction of their total magnetization 
\citep[e.g,][]{li_etal_2021}.


\subsection{Computing the RTP by using a dipole layer}


According to the pioneer work of \citet{silva1986}, which is widely followed in applied geophysics, 
the RTP in the space domain is calculated from total-field anomaly data (Table \partoneref{tbl:Table2}) in two steps.
The first consists in estimating the magnetic-moment intensity distribution on a planar layer 
of dipoles (or doublets) with constant magnetization direction equal to that of the true sources.
The second step consists in computing the $z$-component of the magnetic induction field due 
to the estimated layer by setting its magnetization along the vertical direction.
It is usually verified a posteriori that this computed field approximates the true RTP anomaly (Table \partoneref{tbl:Table2}).

We call attention to the fact that \citet{silva1986} cited the work of \citet{roy1962} to justify that ``the magnetic 
field due to arbitrary sources can be exactly fitted by the magnetic field due to a horizontal equivalent
layer of dipoles''. 
However, \citeauthor{roy1962} neither discussed the existence of a dipole layer that fits the observed total-field 
anomaly data nor its use to retrieve the RTP anomaly, but 
presented a general outline to deduce the solutions of the Dirichlet's and Neumann's problems on a plane 
located above the sources (equation \partoneref{eq:U-Dirichlet-Neumann}).
In other words, the two-step approach proposed by \citet{silva1986}, which is widely accepted by the scientific community 
and we know to be correct from practice 
\citep[e.g., ][]{leao-silva_1989, guspi_novara2009, oliveirajr_etal2013, li_etal2014, reis_etal2020, takahashi-etal2022, piauilino_etal2025}, 
is an empirical generalization of the \citeauthor{roy1962}'s (\citeyear{roy1962}) work. 
It is also important to point out that the mathematical equivalence between the space-domain approach to compute the 
RTP via equivalent-layer technique and the classical approach directly using the Fourier transform of the observed 
total-field anomaly data \citep[e.g., ][ p. 331]{blakely1996} have never been clearly established in the literature.
Here, we provide a general approach aiming at filling these theoretical gaps.

From a theoretical point of view, the two-step approach to compute the RTP in the space domain 
consists in first setting a planar layer of dipoles according to equation \ref{eq:p-Fourier-dipoles}, 
by considering the unit vector $\hat{\mathbf{h}}$ equal to $\hat{\mathbf{v}}$, which defines 
the constant total-magnetization direction of the true sources.
In this case, the ratio $\nicefrac{\left(\mathbf{k}^{\top} \hat{\mathbf{v}} \right)}{\left(\mathbf{k}^{\top} \hat{\mathbf{h}} \right)} = 1$ 
and we can easily verify that the $\mathcal{F} [ p_{\varepsilon}(\mathbf{r}'') ]$ given by 
equation \ref{eq:p-Fourier-dipoles} assumes the form of a Neumann's equivalent layer (Definition \partoneref{def:eql-Neumann}).
The second step to compute the RTP consists in using the predefined Neumann's layer to compute 
the transformed potential field $f_{\varepsilon}$ (equation \partoneref{eq:transformed-field-continuous-layer}) defined in the Fourier domain 
by using $\alpha_{t} = \alpha_{f}$ (equation \partoneref{eq:alpha-t-generic}), with 
$\alpha_{f} \equiv \alpha(2, \hat{\mathbf{z}})$ (equations \partoneref{eq:diff-operators-generic} and \partoneref{eq:canonical-base-vector}).
In this case, it can be verified from equations \partoneref{eq:diff-operators-generic} and \partoneref{eq:transformed-field-continuous-layer-single-step} 
that the computed transformed field $f_{\varepsilon}$ (equation \partoneref{eq:transformed-field-continuous-layer}) 
approximates the RTP anomaly  (Table \partoneref{tbl:Table2}). 
As we have just shown here, for the first time, the classical approach of \citet{silva1986} to compute the RTP in the 
space domain can be formally explained in terms of a Neumann's equivalent layer (Definition \partoneref{def:eql-Neumann}). 

We can also investigate the theoretical link between the RTP in the space domain, 
via equivalent-layer technique, and in the Fourier domain, via classical approach \citep[e.g., ][ p. 331]{blakely1996}. 
As we have already discussed, the transformed field $f_{\varepsilon}$ computed via \partoneref{eq:transformed-field-continuous-layer-Fourier} 
from a previously defined equivalent layer $p_{\varepsilon}$ can also be directly 
computed from the observed potential field $d^{o}$ (equation \partoneref{eq:transformed-field-continuous-layer-1-single-step}
or \partoneref{eq:transformed-field-continuous-layer-2-single-step}), 
in terms of the ratio $\nicefrac{\alpha_{f}}{\alpha_{d}}$.
To compute the RTP, we set $\alpha_{d} \equiv \alpha(1, \hat{\mathbf{u}}) \, \alpha(1, \hat{\mathbf{v}})$ 
and $\alpha_{f} \equiv \alpha(2, \hat{\mathbf{z}})$ (equation \partoneref{eq:diff-operators-generic}).
By using these specific $\alpha_{d}$ and $\alpha_{f}$, setting $z^{\ddagger}_{0} = z^{\dagger}_{0}$ and 
(i) letting $k_{\varepsilon}$ goes to infinity in equation \partoneref{eq:transformed-field-continuous-layer-1-single-step} or
(ii) letting $\mu_{\varepsilon}$ goes to zero in equation \partoneref{eq:transformed-field-continuous-layer-2-single-step}, we get 
\begin{equation}
	\mathcal{F} [ f_{\varepsilon}(\mathbf{r}^{\dagger}) ] =
	\tfrac{k^{2}}{\left(\mathbf{k}^{\top} \hat{\mathbf{u}} \right) \, \left(\mathbf{k}^{\top} \hat{\mathbf{v}} \right)} \, 
	\mathcal{F} \left[ d^{o}(\mathbf{r}^{\dagger}) \right] \: ,
	\label{eq:transformed-field-continuous-layer-single-step-RTP}
\end{equation}
which represents the
classical approach to compute the RTP in the Fourier domain \citep[e.g., ][ p. 331]{blakely1996}.
We stress two important aspects of this result. The first concerns the theoretical link 
between the space- and Fourier-domain approaches to compute the RTP. To the authors knowledge, this link has never been fully 
detailed in the literature.
The second important aspect of our theoretical approach is that it provides the 
reason by which the equivalent-layer technique ultimately suffers the same instability as 
the Fourier-domain RTP operator at low latitudes. In this case, the estimated equivalent layer of dipoles 
presents striations, as have already been pointed out by \citet{li_etal2014}. 
In the Fourier domain, these artifacts are directly related to the term 
$\nicefrac{\left(\mathbf{k}^{\top} \hat{\mathbf{v}} \right)}{\left(\mathbf{k}^{\top} \hat{\mathbf{h}} \right)}$ 
in equation \ref{eq:p-Fourier-dipoles}. This term produces, in the equivalent layer of dipoles, a 
numerical instability similar to that caused by the term 
$\nicefrac{k^{2}}{\left(\mathbf{k}^{\top} \hat{\mathbf{u}} \right) \, \left(\mathbf{k}^{\top} \hat{\mathbf{v}} \right)}$ in equation 
\ref{eq:transformed-field-continuous-layer-single-step-RTP},
which is the base of the classical approach to compute the RTP in the Fourier domain \citep[e.g., ][ p. 331]{blakely1996}.


\subsection{Generalized equivalent sources}

Let us define an equivalent layer $p_{\varepsilon}$ by considering the kernel function 
$a(\mathbf{r}, \mathbf{r}'') = \Psi(\mathbf{r}, \mathbf{r}'')$ (equation \partoneref{eq:Psi}). 
Note that, in this case, 
the corresponding predicted potential field $d_{\varepsilon}$ (equation \partoneref{eq:predicted-data-continuous-layer}) 
is analogous to the gravitational potential due to the equivalent layer.
Besides, according to equations \partoneref{eq:p-Fourier}--\partoneref{eq:kernel-a-Fourier}, we have that 
$\alpha_{a} = 1$ and 
\begin{equation}
	\mathcal{F} [ p_{\varepsilon}(\mathbf{r}'') ] =
	k \, \alpha_{d} \, \mathcal{F} \left[ U_{\varepsilon}(\mathbf{r}'') \right]  \: ,
	\label{eq:p-Fourier-generalized-eqs}
\end{equation}
which approximates the vertical derivative of the observed potential field on the layer plane. 
Such layers produce potential-field transformations (equations \partoneref{eq:transformed-field-continuous-layer} and 
\partoneref{eq:transformed-field-continuous-layer-Fourier}) defined by $\alpha_{t} = \nicefrac{\alpha_{f}}{\alpha_{d}}$ (equation \partoneref{eq:alpha-t-generic}).

This strategy, called \textit{generalized equivalent sources}, was proposed by \citet{cordell1992} and recently adapted by 
\citet{soler_uieda2021} to perform gravity-data processing. 
Those authors considered only the particular case in which the observed potential field $d^{o}$ is gravity disturbance data 
$\tfrac{\partial}{\partial z} U(\mathbf{r})$  (Table \partoneref{tbl:Table2}), 
with $\alpha_{d} \equiv \alpha(1, \hat{\mathbf{z}}) = k$ (equations \partoneref{eq:observed-data-Fourier}, 
\partoneref{eq:diff-operators-generic} and \partoneref{eq:canonical-base-vector}).
In this particular case, equation \ref{eq:p-Fourier-generalized-eqs} leads to an equivalent layer that 
(i) approximates $\tfrac{\partial^{2}}{\partial z^{2}} U(\mathbf{r})$ (equation \partoneref{eq:potential-U-2nd-deriv-space}) 
in the space domain, 
(ii) has minimum phase changes (Definition \partoneref{def:eql-min-phase}) and 
(iii) results in potential-field transformations defined by $\alpha_{t} = \nicefrac{\alpha_{f}}{k}$ (equation \partoneref{eq:alpha-t-generic}).

\subsection{Alternative approach to compute the RTP}

\citet{guspi_novara2009} proposed an alternative two-step approach to compute the RTP via equivalent-layer technique. 
The first step consists in obtaining an equivalent layer by considering that (i) the observed potential field $d^{o}$
is satisfactorily described by the directional derivative 
$\hat{\mathbf{u}}^{\top} \mathbf{H}_{U}(\mathbf{r}) \, \hat{\mathbf{v}}$ 
(equation \partoneref{eq:potential-U-oblique-2nd-deriv-space}), where $\hat{\mathbf{u}}$ defines the constant direction of the main geomagnetic
field at the study area and $\hat{\mathbf{v}}$ the constant total-magnetization direction of the sources  (Table \partoneref{tbl:Table2}), and (ii)
the kernel function is given by 
\begin{equation}
	a(\mathbf{r}, \mathbf{r}'') = \tfrac{\left( z''_{0} - z \right)}{2 \, \pi} \, 
	\ln{\Big[} \left( z''_{0} - z \right) + \| \mathbf{r}'' - \mathbf{r} \| {\Big]} - 
	\tfrac{\| \mathbf{r}'' - \mathbf{r} \|}{2 \, \pi} \: ,
	\label{eq:kernel-Fourier-Guspi}
\end{equation}
which represents the double integral of the function $\Psi$ (equation \partoneref{eq:Psi}) 
along the vertical direction $z$.
The second step of the alternative approach proposed by \citet{guspi_novara2009} computes the RTP anomaly 
(Table \partoneref{tbl:Table2}) by setting $f_{\varepsilon}$ (equation \partoneref{eq:transformed-field-continuous-layer}) with 
$t(\mathbf{r}, \mathbf{r}'') = \Psi(\mathbf{r}, \mathbf{r}'')$ (equation \partoneref{eq:Psi}).

Note that the first step consists in obtaining an equivalent layer $p_{\varepsilon}$ defined by 
$\alpha_{d} = (\mathbf{k}^{\top} \hat{\mathbf{u}}) (\mathbf{k}^{\top} \hat{\mathbf{v}})$ (equations \partoneref{eq:observed-data-Fourier} 
and \partoneref{eq:diff-operators-generic}) and 
$\alpha_{a} = \nicefrac{(\mathbf{k}^{\top} \hat{\mathbf{u}}) (\mathbf{k}^{\top} \hat{\mathbf{h}})}{k^{2}}$, 
where the term $\nicefrac{1}{k^{2}}$ represents the double integration along the vertical direction. In this case,  
$p_{\varepsilon}$ assumes the following form in the Fourier domain:
\begin{equation}
	\mathcal{F} [ p_{\varepsilon}(\mathbf{r}'') ] =
	k^{3} \, \tfrac{\left(\mathbf{k}^{\top} \hat{\mathbf{v}} \right)}{\left(\mathbf{k}^{\top} \hat{\mathbf{h}} \right)} \,
	\mathcal{F} \left[ U_{\varepsilon}(\mathbf{r}'') \right] \: .
	\label{eq:p-Fourier-guspi}
\end{equation}
The second step consists in computing $f_{\varepsilon}$ in the Fourier domain (equation \partoneref{eq:transformed-field-continuous-layer-Fourier}) 
by using $\alpha_{t} = 1$ (equation \partoneref{eq:alpha-t-generic}).
In this case, it can be verified that the transformed potential field $\mathcal{F} [ f_{\varepsilon}(\mathbf{r}^{\ddagger}) ]$ 
(equation \partoneref{eq:transformed-field-continuous-layer-Fourier}) actually approximates the RTP anomaly  (Table \partoneref{tbl:Table2}) 
for the particular case in which $\hat{\mathbf{h}} = \hat{\mathbf{v}}$, 
i.e., for the case in which the constant magnetization direction of the equivalent layer coincides with that of the true sources.
This result shows the effectiveness of our approach in providing the theoretical background not only for classical 
equivalent-layer methods, but also for unconventional ones.